% =========================================================================
% File        : 06_plotting.tex
% Deskripsi   : Slide "Membuat Plot di MATLAB" (Basic MATLAB, Modul 06).
%               Rekonstruksi dari PDF terenkripsi materi/06.bin; gambar
%               figs/f01--f28 dihasilkan oleh buat_gambar_plot.m.
% Proyek      : Basic MATLAB -- Modul Plotting
% Penulis     : Andreas Syaloom Kurniawan
% Versi       : 1.1 (rekonstruksi)
% Dibuat      : 13 Oktober 2026
% Diperbarui  : 14 Oktober 2026
% =========================================================================
% Kompilasi: xelatex 06_plotting.tex  (butuh font Fira Sans & Fira Mono)

\documentclass[10pt,aspectratio=169]{beamer}
\usetheme[numbering=counter]{metropolis}

\usepackage{fontspec}
\setsansfont{Fira Sans}[BoldFont={Fira Sans Bold},ItalicFont={Fira Sans Italic},BoldItalicFont={Fira Sans Bold Italic}]
\setmonofont{Fira Mono}[BoldFont={Fira Mono Bold}]

\usepackage{booktabs}
\usepackage{array}
\usepackage{listings}
\usepackage{tikz}
\usetikzlibrary{arrows.meta,positioning}
\usepackage{pgfplots}
\pgfplotsset{compat=1.17}

\setbeamerfont{title}{size=\Large,series=\bfseries}
\setbeamerfont{subtitle}{size=\large}
\setbeamerfont{frametitle}{size=\large,series=\bfseries}
\setbeamerfont{section title}{size=\Large,series=\bfseries}
\setbeamerfont{institute}{size=\small}
\setbeamerfont{page number in head/foot}{size=\scriptsize}

% ---- Warna --------------------------------------------------------------
\definecolor{mblue}{HTML}{0072BD}
\definecolor{morange}{HTML}{D95319}
\definecolor{kodebg}{HTML}{F5F5F5}
\definecolor{kodekw}{HTML}{0000CC}
\definecolor{kodestr}{HTML}{A020F0}
\definecolor{kodecom}{HTML}{228B22}

% ---- Listing MATLAB -----------------------------------------------------
\lstdefinestyle{mat}{
  language=Matlab,
  basicstyle=\ttfamily\scriptsize,
  keywordstyle=\bfseries\color{kodekw},
  stringstyle=\color{kodestr},
  commentstyle=\color{kodecom},
  morekeywords={on,off,show,minor,left,right,histogram,xline,yline,
                yyaxis,tiledlayout,nexttile,exportgraphics},
  backgroundcolor=\color{kodebg},
  frame=single, rulecolor=\color{black!25}, framerule=0.4pt,
  framesep=3pt, xleftmargin=3pt, xrightmargin=3pt,
  aboveskip=4pt, belowskip=4pt,
  columns=fullflexible, keepspaces=true,
  showstringspaces=false, upquote=true
}
\lstset{style=mat}

% ---- Makro --------------------------------------------------------------
\newcommand{\konsep}{\alert{\textbf{Konsep.}}\ }
\newcommand{\kd}[1]{\texttt{#1}}
% Tata letak "teks kiri, gambar kanan": \kiri{lebar gambar} ... \kanan{lebar}{file}
\newcommand{\kiri}[1]{\begin{columns}[T,onlytextwidth]
  \column{\dimexpr\textwidth-#1-1.2em\relax}\small}
\newcommand{\kanan}[2]{\column{#1}\includegraphics[width=#1]{figs/#2}\end{columns}}

\title{Membuat Plot di MATLAB}
\subtitle{Dari data hingga grafik yang siap dipresentasikan}
\institute{Pemrograman MATLAB}
\date{}

\begin{document}

\maketitle
\setcounter{framenumber}{0}

% =========================================================================
\begin{frame}{Tujuan pembelajaran}
Setelah mempelajari materi ini, peserta diharapkan mampu:
\begin{enumerate}
  \item Menjelaskan \alert{anatomi} grafik MATLAB: \emph{figure}, \emph{axes}, dan objek grafik.
  \item Menyiapkan data untuk plot, baik dari vektor maupun dari objek \kd{table}.
  \item \alert{Memilih} jenis grafik yang sesuai dengan pertanyaan yang ingin dijawab (garis, batang, sebar, dan lainnya).
  \item Mengatur tampilan grafik: judul, label, legenda, warna, gaya garis, marker, dan sumbu.
  \item Menggabungkan beberapa grafik dalam satu figure lalu \alert{menyimpannya} untuk laporan atau slide.
\end{enumerate}

\medskip
{\small Prasyarat: membuat vektor (\kd{1:5}, \kd{linspace}), operasi elemen demi elemen (\kd{.*}, \kd{.\textasciicircum}), dan objek \kd{table}.}
\end{frame}

\begin{frame}{Isi presentasi}
\begin{columns}[c,onlytextwidth]
\column{0.62\textwidth}
\begin{enumerate}
  \item Konsep dasar grafik
  \item Sumber data
  \item Plot dasar dan memilih jenis grafik
  \item Teks dan keterangan
  \item Gaya garis, marker, dan sumbu
  \item Beberapa grafik dan menyimpan hasil
  \item Penutup
\end{enumerate}
\column{0.33\textwidth}\small
\textbf{Pola tiap slide}\par\smallskip
\alert{\textbf{Konsep.}} ide singkat\\
$\downarrow$\\
Kode MATLAB\\
$\downarrow$\\
Catatan \& hasil grafik\par\medskip
Ketik ulang kodenya sendiri di MATLAB, jangan hanya dibaca.
\end{columns}
\end{frame}

\begin{frame}[fragile]{Data yang dipakai sepanjang materi}
Semua contoh memakai data ilustrasi tiga indikator makroekonomi berikut (\alert{bukan} data resmi). Jalankan kode ini sekali di awal agar semua contoh bisa langsung dicoba.
\begin{lstlisting}[basicstyle=\ttfamily\footnotesize]
tahun   = 2018:2024;                         % 7 tahun
growth  = [5.17 5.02 -2.07 3.70 5.31 5.05 5.03];  % pertumbuhan ekonomi (%)
inflasi = [3.13 2.72  1.68 1.87 5.51 2.61 1.57];  % inflasi (%)
bunga   = [6.00 5.00  3.75 3.50 5.75 6.25 6.00];  % suku bunga (%)
\end{lstlisting}
\begin{center}\small
\begin{tabular}{lrrrrrrr}
\toprule
Tahun & 2018 & 2019 & 2020 & 2021 & 2022 & 2023 & 2024\\
\midrule
Growth (\%)     & 5,17 & 5,02 & $-$2,07 & 3,70 & 5,31 & 5,05 & 5,03\\
Inflasi (\%)    & 3,13 & 2,72 & 1,68 & 1,87 & 5,51 & 2,61 & 1,57\\
Suku bunga (\%) & 6,00 & 5,00 & 3,75 & 3,50 & 5,75 & 6,25 & 6,00\\
\bottomrule
\end{tabular}
\end{center}
{\small Keempat vektor berukuran $1\times 7$. Ingat: MATLAB memakai \alert{titik} sebagai pemisah desimal (\kd{5.17}), bukan koma.}
\end{frame}

% =========================================================================
\section{Konsep dasar grafik}

\begin{frame}{Anatomi sebuah grafik MATLAB}
\begin{columns}[T,onlytextwidth]
\column{0.44\textwidth}
\begin{tikzpicture}
  \node[draw=black!70, line width=0.6pt, fill=white, inner sep=4pt, anchor=north west] (fig) at (0,0) {%
  \begin{tikzpicture}
  \begin{axis}[
    width=6.4cm, height=4.7cm,
    title={\textbf{Indikator Ekonomi}}, title style={font=\scriptsize},
    xlabel={Tahun}, ylabel={Persen (\%)},
    label style={font=\scriptsize}, tick label style={font=\scriptsize},
    xtick={2018,2020,2022}, ytick={0,4},
    x tick label style={/pgf/number format/1000 sep={}},
    legend style={font=\scriptsize, at={(0.97,0.03)}, anchor=south east},
    legend cell align=left, every axis plot/.append style={thick}]
    \addplot[color=mblue, mark=none] coordinates
      {(2018,5.17)(2019,5.02)(2020,-2.07)(2021,3.70)(2022,5.31)(2023,5.05)(2024,5.03)};
    \addplot[color=morange, mark=none] coordinates
      {(2018,3.13)(2019,2.72)(2020,1.68)(2021,1.87)(2022,5.51)(2023,2.61)(2024,1.57)};
    \legend{Growth,Inflasi}
  \end{axis}
  \end{tikzpicture}};
  \node[anchor=north west, font=\scriptsize\bfseries] at (fig.north west) {Figure 1};
\end{tikzpicture}
\column{0.52\textwidth}\small
\textbf{Tiga lapis objek:}
\begin{enumerate}
  \item \textbf{Figure}: jendela tempat grafik digambar. Dibuka dengan \kd{figure}.
  \item \textbf{Axes}: bidang koordinat di dalam figure (sumbu, \emph{tick}, grid). Satu figure bisa berisi beberapa axes.
  \item \textbf{Objek grafik}: garis (\kd{plot}), batang (\kd{bar}), titik (\kd{scatter}), dst.
\end{enumerate}
Judul (\kd{title}), label sumbu (\kd{xlabel}, \kd{ylabel}), dan legenda (\kd{legend}) adalah \alert{keterangan} yang menempel pada axes.
\end{columns}
\end{frame}

\begin{frame}{Alur kerja membuat grafik}
\begin{center}
\begin{tikzpicture}[
  langkah/.style={draw=black!25, fill=black!4, rounded corners=2pt,
    align=center, font=\footnotesize, minimum height=2.3em, inner xsep=5pt},
  panah/.style={-{Latex[length=4pt]}, thick, color=mLightBrown}]
  \node[langkah] (a) {\textbf{1. Siapkan}\\data};
  \node[langkah, right=0.5cm of a] (b) {\textbf{2. Buka}\\\texttt{figure}};
  \node[langkah, right=0.5cm of b] (c) {\textbf{3. Gambar}\\\texttt{plot, bar, \ldots}};
  \node[langkah, right=0.5cm of c] (d) {\textbf{4. Lengkapi}\\judul, label, legenda};
  \node[langkah, right=0.5cm of d] (e) {\textbf{5. Simpan}\\\texttt{exportgraphics}};
  \draw[panah] (a) -- (b); \draw[panah] (b) -- (c);
  \draw[panah] (c) -- (d); \draw[panah] (d) -- (e);
\end{tikzpicture}
\end{center}
\begin{itemize}
  \item Langkah 1 menentukan \alert{apa} yang digambar; langkah 3 menentukan \alert{bagaimana} data ditampilkan.
  \item Langkah 4 sering dilewatkan, padahal grafik tanpa label dan satuan tidak bisa dibaca orang lain.
  \item Untuk menambah objek ke axes yang sama (langkah 3 berulang), pakai \kd{hold on}.
\end{itemize}
\smallskip
{\small Prinsip: grafik yang baik bisa dipahami \alert{tanpa} penjelasan lisan dari pembuatnya.}
\end{frame}

\begin{frame}[fragile]{Pola umum pemanggilan fungsi plot}
\small
\konsep Hampir semua fungsi grafik mengikuti pola yang sama: data dulu, lalu pilihan gaya, lalu pengaturan detail.
\begin{lstlisting}[basicstyle=\ttfamily\footnotesize]
plot(x, y, LineSpec, 'NamaProperti', nilai, 'NamaProperti', nilai, ...)
%    ---- data  --------  ---------------- pasangan nama-nilai -------------
\end{lstlisting}
\begin{tabular}{p{0.2\textwidth}p{0.38\textwidth}l}
\toprule
Bagian & Fungsi & Contoh\\
\midrule
Data & Wajib. Koordinat yang digambar & \kd{plot(tahun, growth)}\\
\emph{LineSpec} & Opsional. Singkatan warna, gaya garis, marker dalam satu teks & \kd{'r--o'}\\
Pasangan nama--nilai & Opsional. Mengatur properti tertentu secara rinci & \kd{'LineWidth', 2}\\
\bottomrule
\end{tabular}

\medskip
Contoh: \kd{plot(tahun, growth, 'r--o', 'LineWidth', 2)}. Sejak R2021a boleh juga \kd{LineWidth=2}. Nama properti tidak peka huruf besar/kecil.
\end{frame}

% =========================================================================
\section{Sumber data}

\begin{frame}[fragile]{Sumber data untuk plotting}
\small
\konsep Plot hanya menerima angka dalam bentuk \alert{vektor} atau matriks. Dari mana pun asal datanya, ujungnya harus menjadi vektor.
\begin{columns}[T,onlytextwidth]
\column{0.48\textwidth}
\textbf{1. Diketik langsung}
\begin{lstlisting}
x = 1:5;
y = [2 4 3 8 5];
t = linspace(0, 2*pi, 100);
y2 = sin(t);
\end{lstlisting}
\textbf{2. File CSV / Excel}
\begin{lstlisting}
D = readtable('ekonomi.csv');
D = readtable('data.xlsx');
tahun  = D.Tahun;
growth = D.Growth;
\end{lstlisting}
\column{0.48\textwidth}
\textbf{3. File MATLAB (.mat)}
\begin{lstlisting}
load('data.mat')
% isi file jadi variabel
\end{lstlisting}
\textbf{4. Hasil perhitungan}
\begin{lstlisting}
x = 0:0.1:10;
y = x.^2 - 3*x;  % .^ per elemen
\end{lstlisting}
Syarat utama: \kd{x} dan \kd{y} punya \alert{jumlah elemen yang sama}. Cek dengan \kd{numel(x)} dan \kd{numel(y)}. Vektor baris dan kolom boleh dicampur selama panjangnya sama.
\end{columns}
\end{frame}

\begin{frame}[fragile]{Sumber data: objek tabel (\texttt{table})}
\kiri{198pt}
\konsep Data dari file biasanya masuk sebagai \kd{table}: setiap kolom punya nama. Satu kolom tabel = satu vektor.
\begin{lstlisting}
D = readtable('ekonomi.csv');

class(D)                   % 'table'
D.Properties.VariableNames
% Tahun Growth Inflasi SukuBunga

plot(D.Tahun, D.Growth, '-o')
xlabel('Tahun')
ylabel('Growth (%)')
\end{lstlisting}
Kolom diakses dengan \kd{NamaTabel.NamaKolom}; hasilnya vektor kolom biasa sehingga bisa langsung masuk ke \kd{plot}. Cara lain: \kd{D\{:,'Growth'\}}.
\kanan{198pt}{f18_tabel1}
\end{frame}

\begin{frame}[fragile]{Plot langsung dari tabel}
\kiri{198pt}
\konsep Tabel bisa diberikan langsung ke \kd{plot}; cukup sebut nama kolom untuk sumbu $x$ dan $y$.
\begin{lstlisting}
% x = kolom Tahun,
% y = tiga kolom sekaligus
plot(D, 'Tahun', ...
     {'Growth','Inflasi','SukuBunga'})
xlabel('Tahun')
ylabel('Persen (%)')
legend('Location','southoutside', ...
       'Orientation','horizontal')
\end{lstlisting}
Bentuk \kd{plot(tabel, 'kolomX', 'kolomY')} ada sejak R2022a. Beberapa kolom $y$ ditulis dalam \kd{\{ \}}; nama kolom otomatis menjadi nama garis di legenda. Versi lama: pakai \kd{D.Kolom}.
\kanan{198pt}{f19_tabel2}
\end{frame}

% =========================================================================
\section{Plot dasar dan memilih jenis grafik}

\begin{frame}[fragile]{Plot sederhana: hanya $y$}
\kiri{207pt}
\konsep Jika hanya diberi satu vektor, MATLAB tetap bisa menggambar, tetapi tidak tahu ``waktunya''.
\begin{lstlisting}
plot(growth)
\end{lstlisting}
Sumbu $x$ otomatis berupa \alert{indeks} $1, 2, \ldots, 7$ (urutan elemen), bukan tahun. Cocok untuk melihat bentuk data sekilas, tetapi \alert{tidak} untuk disajikan.
\kanan{207pt}{f01_sederhana}
\end{frame}

\begin{frame}[fragile]{Plot grafik $(x, y)$}
\kiri{207pt}
\konsep Berikan pasangan $(x, y)$: elemen ke-$i$ dari \kd{x} dipasangkan dengan elemen ke-$i$ dari \kd{y}, lalu titik-titiknya dihubungkan garis.
\begin{lstlisting}
plot(tahun, growth)
xlabel('Tahun')
ylabel('Pertumbuhan ekonomi (%)')
\end{lstlisting}
\kd{x} di sumbu mendatar, \kd{y} di sumbu tegak. Sekarang sumbu $x$ sudah berisi tahun. \kd{xlabel} dan \kd{ylabel} memberi nama sumbu, \alert{sertakan satuannya}.
\kanan{207pt}{f02_xy}
\end{frame}

\begin{frame}[fragile]{Beberapa garis dalam satu plot}
\kiri{207pt}
\konsep Tulis beberapa pasangan $(x, y)$ berurutan dalam satu perintah \kd{plot}; setiap pasangan menjadi satu garis.
\begin{lstlisting}
plot(tahun, growth, ...
     tahun, inflasi, ...
     tahun, bunga)
xlabel('Tahun')
ylabel('Persen (%)')
legend('Growth','Inflasi', ...
       'Suku bunga')
\end{lstlisting}
Jumlah garis tidak dibatasi. Warna diambil otomatis dari urutan warna bawaan (biru, oranye, kuning, \ldots). Karena ada lebih dari satu garis, \alert{legenda wajib} ada.
\kanan{207pt}{f03_multi}
\end{frame}

\begin{frame}[fragile]{Cara lain: plot sebuah matriks}
\kiri{198pt}
\konsep Jika \kd{y} berupa matriks, MATLAB menggambar \alert{satu garis untuk setiap kolom}.
\begin{lstlisting}
Y = [growth; inflasi; bunga]';
size(Y)      % 7 3

plot(tahun, Y)
xlabel('Tahun')
ylabel('Persen (%)')
legend('Growth','Inflasi', ...
       'Suku bunga')
\end{lstlisting}
Ketiga vektor baris ditumpuk menjadi matriks $3 \times 7$, lalu ditranspos (\kd{'}) menjadi $7 \times 3$ agar setiap \alert{kolom} adalah satu variabel. Hasilnya sama dengan slide sebelumnya, tetapi kodenya lebih ringkas jika variabelnya banyak.
\kanan{198pt}{f03_multi}
\end{frame}

\begin{frame}[fragile]{Mengelola jendela: \texttt{figure} dan \texttt{hold}}
\small
\konsep Setiap \kd{plot} baru \alert{menghapus} isi axes yang sedang aktif. \kd{hold on} ``menahan'' isi lama sehingga objek berikutnya ditambahkan.
\begin{columns}[T,onlytextwidth]
\column{0.42\textwidth}
\begin{lstlisting}
figure                 % jendela baru
plot(tahun, growth)
plot(tahun, inflasi)   % menimpa!

figure
plot(tahun, growth)
hold on                % tahan isi
plot(tahun, inflasi)   % ditambah
hold off
\end{lstlisting}
\begin{tabular}{@{}ll@{}}
\kd{figure}    & buka jendela baru\\
\kd{clf}       & kosongkan figure aktif\\
\kd{close all} & tutup semua jendela\\
\end{tabular}
\column{222pt}
\includegraphics[width=222pt]{figs/f21_hold}\par\medskip
Biasakan memanggil \kd{figure} sebelum grafik baru agar grafik lama tidak tertimpa.
\end{columns}
\end{frame}

\begin{frame}[fragile]{Diagram batang}
\kiri{207pt}
\konsep Batang menonjolkan \alert{besar nilai} tiap kategori atau periode, bukan arah perubahannya.
\begin{lstlisting}
bar(tahun, growth)
xlabel('Tahun')
ylabel('Pertumbuhan ekonomi (%)')
\end{lstlisting}
Tinggi batang = nilai data. Nilai negatif tampak jelas di bawah garis nol (kontraksi 2020). Lebar batang bisa diatur dengan argumen ketiga, misalnya \kd{bar(tahun, growth, 0.5)}.
\kanan{207pt}{f04_bar}
\end{frame}

\begin{frame}[fragile]{Diagram batang berkelompok}
\kiri{207pt}
\konsep Sama seperti \kd{plot}, jika \kd{bar} diberi matriks, \alert{setiap kolom} menjadi satu kelompok warna.
\begin{lstlisting}
B = [growth; inflasi]';   % 7x2
bar(tahun, B)
xlabel('Tahun')
ylabel('Persen (%)')
legend('Growth','Inflasi', ...
       'Location','southwest')
\end{lstlisting}
Setiap baris matriks = satu tahun (satu kelompok batang), setiap kolom = satu variabel. Variasi \kd{bar(tahun, B, 'stacked')} menumpuk batang, cocok untuk komposisi.
\kanan{207pt}{f22_bar_grup}
\end{frame}

\begin{frame}[fragile]{Scatterplot}
\kiri{207pt}
\konsep Diagram sebar menampilkan \alert{hubungan} dua variabel: setiap titik adalah satu pengamatan (di sini: satu tahun).
\begin{lstlisting}
scatter(inflasi, growth, 60, 'filled')
xlabel('Inflasi (%)')
ylabel('Pertumbuhan ekonomi (%)')
\end{lstlisting}
Pola umumnya \kd{scatter(x, y, ukuran)}. Titik tidak dihubungkan garis karena urutan data tidak bermakna. \kd{'filled'} mengisi titik dengan warna.
\kanan{207pt}{f05_scatter}
\end{frame}

\begin{frame}[fragile]{Scatterplot dengan garis tren}
\kiri{198pt}
\konsep Garis tren $y = a + b\,x$ membantu membaca arah hubungan. Koefisien dihitung dengan \kd{polyfit}, lalu digambar di atas titik.
\begin{lstlisting}
scatter(inflasi, growth, 60, 'filled')
hold on
p  = polyfit(inflasi, growth, 1);  % [b a]
xs = linspace(min(inflasi), max(inflasi));
plot(xs, polyval(p, xs), 'r-', ...
     'LineWidth', 1.5)
hold off
xlabel('Inflasi (%)')
ylabel('Pertumbuhan ekonomi (%)')
\end{lstlisting}
\kd{polyfit(x,y,1)}: garis lurus terbaik (kuadrat terkecil); \kd{polyval}: nilainya di titik \kd{xs}. Dengan 7 titik dan satu pencilan (2020), garis ini hanya ilustrasi, \alert{bukan} bukti sebab-akibat.
\kanan{198pt}{f23_scatter_tren}
\end{frame}

\begin{frame}[fragile]{Jenis grafik lain yang sering dipakai}
\kiri{198pt}
\begin{lstlisting}
% sebaran (distribusi) data
histogram(randn(500,1) + 5, 20)

% luas bertumpuk (komposisi)
area(tahun, [inflasi; bunga]')

% tangga: nilai tetap sampai berubah
stairs(tahun, bunga)

% proporsi dari keseluruhan
pie([45 30 25])
\end{lstlisting}
Pola pemanggilannya sama: data dulu, lalu pilihan. \kd{stairs} cocok untuk suku bunga kebijakan yang berubah bertahap. \kd{pie} sebaiknya dipakai hemat: mata manusia sulit membandingkan sudut.
\kanan{198pt}{f24_jenis_lain}
\end{frame}

\begin{frame}{Memilih jenis grafik: mulai dari pertanyaan}
\small
\konsep Jenis grafik ditentukan oleh \alert{pertanyaan} yang ingin dijawab, bukan oleh selera.

\medskip
\begin{tabular}{p{0.45\textwidth}p{0.2\textwidth}p{0.25\textwidth}}
\toprule
Pertanyaan & Grafik & Perintah\\
\midrule
Bagaimana nilai berubah dari waktu ke waktu? & Garis & \kd{plot}\\
Berapa besar nilai tiap kategori/periode? & Batang & \kd{bar}\\
Apakah dua variabel saling berhubungan? & Sebar & \kd{scatter}\\
Bagaimana sebaran (distribusi) data? & Histogram & \kd{histogram}\\
Bagaimana komposisi suatu total? & Batang bertumpuk & \kd{bar(\ldots,'stacked')}, \kd{area}\\
Dua variabel dengan skala/satuan berbeda? & Dua sumbu $y$ & \kd{yyaxis}\\
Data tumbuh eksponensial / rentang sangat lebar? & Skala log & \kd{semilogy}\\
\bottomrule
\end{tabular}
\end{frame}

% =========================================================================
\section{Teks dan keterangan}

\begin{frame}[fragile]{Judul dan subjudul}
\kiri{207pt}
\konsep Judul menjawab ``grafik ini tentang apa''; subjudul memberi konteks (sumber, periode).
\begin{lstlisting}
plot(tahun, growth)
title('Pertumbuhan Ekonomi')
subtitle('Data ilustrasi, 2018-2024')
xlabel('Tahun')
ylabel('Persen (%)')
\end{lstlisting}
\kd{subtitle} tersedia sejak R2020b. Pada versi lama: \kd{title(\{'baris 1','baris 2'\})} untuk judul dua baris. Untuk grafik di makalah/slide, judul sering dihilangkan karena sudah ada di \emph{caption}: cukup \alert{tidak memanggil} \kd{title}.
\kanan{207pt}{f06_judul}
\end{frame}

\begin{frame}[fragile]{Karakter khusus pada teks: \emph{interpreter}}
\small
\konsep Teks pada judul, label, dan legenda dibaca dengan aturan \TeX: \kd{\_} membuat subskrip, \kd{\textasciicircum} membuat superskrip. Ini sering membuat nama variabel tampil salah.
\begin{columns}[T,onlytextwidth]
\column{0.48\textwidth}
\begin{lstlisting}
title('GDP_growth')
% tampil: GDP dengan subskrip "g"

title('GDP_growth', 'Interpreter','none')
title('GDP\_growth')   % atau escape
\end{lstlisting}
Hasil tanpa pengaturan: GDP$_\mathrm{g}$rowth. Dengan \kd{'Interpreter','none'} atau \kd{\textbackslash\_}: GDP\_growth.
\column{0.48\textwidth}
\begin{lstlisting}
% huruf Yunani
xlabel('\alpha dan \beta')
% pangkat dan indeks
ylabel('x^2 + y_t')
% rumus LaTeX
title('$\hat{y}_t$', ...
      'Interpreter','latex')
\end{lstlisting}
Interpreter bawaan \kd{'tex'} sudah mendukung huruf Yunani dan pangkat sederhana. Untuk rumus yang lebih rumit, pakai \kd{'latex'} dan apit rumus dengan \$\ldots\$.
\end{columns}
\end{frame}

\begin{frame}[fragile]{Membuat legenda}
\kiri{182pt}
\konsep Legenda menghubungkan warna/gaya garis dengan nama variabel. Ada dua cara memberi nama.
\begin{lstlisting}
% Cara 1: beri nama saat menggambar
hold on
plot(tahun, growth,  'DisplayName','Growth')
plot(tahun, inflasi, 'DisplayName','Inflasi')
plot(tahun, bunga,   'DisplayName','Suku bunga')
hold off
legend show

% Cara 2: beri nama sesudahnya
% legend('Growth','Inflasi','Suku bunga')
\end{lstlisting}
Cara 2 mengikuti \alert{urutan} garis digambar; salah urutan berarti salah label. Cara 1 lebih aman karena nama melekat pada garisnya.
\kanan{182pt}{f08_legend}
\end{frame}

\begin{frame}[fragile]{Memosisikan legenda}
\kiri{182pt}
\konsep Legenda tidak boleh menutupi data. Atur posisinya dengan \kd{'Location'} dan arah susunannya dengan \kd{'Orientation'}.
\begin{lstlisting}
plot(tahun, growth, ...
     tahun, inflasi, ...
     tahun, bunga)
legend('Growth','Inflasi','Suku bunga', ...
  'Location','southoutside', ...
  'Orientation','horizontal')
\end{lstlisting}
Arah mata angin: \kd{'north'} (atas), \kd{'south'} (bawah), \kd{'east'} (kanan), \kd{'west'} (kiri), serta gabungannya (\kd{'northwest'}, \ldots). Akhiran \kd{outside} menaruh legenda \alert{di luar} axes. \kd{'best'} membiarkan MATLAB mencari tempat yang paling sedikit menutupi data.
\kanan{182pt}{f13_posisi}
\end{frame}

\begin{frame}[fragile]{Anotasi: garis acuan dan teks}
\kiri{182pt}
\konsep Anotasi mengarahkan perhatian pembaca ke peristiwa penting, misalnya batas nol atau tahun krisis.
\begin{lstlisting}
plot(tahun, growth, '-o')
yline(0, 'k-')       % garis y = 0
xline(2020, '--r', 'Pandemi COVID-19')
text(2021.1, 3.7, '\leftarrow pemulihan')
xlabel('Tahun')
ylabel('Growth (%)')
\end{lstlisting}
\kd{xline}/\kd{yline} (sejak R2018b) menggambar garis tegak/mendatar sepanjang axes, dengan label opsional. \kd{text(x, y, 'teks')} menaruh tulisan pada \alert{koordinat data}. \kd{\textbackslash leftarrow} adalah panah dari interpreter \TeX.
\kanan{182pt}{f27_anotasi}
\end{frame}

% =========================================================================
\section{Gaya garis, marker, dan sumbu}

\begin{frame}[fragile]{Tebal garis}
\kiri{182pt}
\konsep Ketebalan membedakan garis utama dari garis pendukung. Diatur dengan properti \kd{'LineWidth'}.
\begin{lstlisting}
hold on
plot(tahun, growth,  'LineWidth',0.5)
plot(tahun, inflasi, 'LineWidth',2)
plot(tahun, bunga,   'LineWidth',4)
hold off
\end{lstlisting}
Satuannya \emph{point} (1 pt $\approx$ 0,35 mm); bawaan MATLAB 0,5. Untuk grafik di slide atau cetakan, 1,5--2 biasanya mudah dibaca.
\kanan{182pt}{f09_tebal}
\end{frame}

\begin{frame}[fragile]{Warna garis}
\kiri{182pt}
\konsep Warna bisa ditentukan dengan tiga cara, dari yang paling singkat sampai yang paling tepat.
\begin{lstlisting}
hold on
plot(tahun, growth,  'r')                % 1
plot(tahun, inflasi, 'Color',[0 0.5 0])  % 2
plot(tahun, bunga,   'Color','#FF8800')  % 3
hold off
\end{lstlisting}
(1) Huruf singkat hanya 8 warna dasar. (2) Vektor \kd{[r g b]} dengan nilai \alert{0--1}, bukan 0--255: \kd{[0 0.5 0]} = hijau tua. (3) Kode heksadesimal \kd{'\#RRGGBB'} (sejak R2019a), praktis jika menyalin warna dari palet.
\kanan{182pt}{f10_warna}
\end{frame}

\begin{frame}{Referensi singkat: \emph{LineSpec}}
\small
\begin{columns}[T,onlytextwidth]
\column{0.3\textwidth}
\begin{tabular}{@{}cl@{}}
\toprule
\multicolumn{2}{c}{Warna}\\
\midrule
\kd{r} & merah\\ \kd{g} & hijau\\ \kd{b} & biru\\ \kd{k} & hitam\\
\kd{y} & kuning\\ \kd{m} & magenta\\ \kd{c} & cyan\\ \kd{w} & putih\\
\bottomrule
\end{tabular}
\column{0.3\textwidth}
\begin{tabular}{@{}cl@{}}
\toprule
\multicolumn{2}{c}{Gaya garis}\\
\midrule
\kd{-}  & lurus (bawaan)\\
\kd{-{}-} & putus-putus\\
\kd{:}  & titik-titik\\
\kd{-.} & garis-titik\\
\bottomrule
\end{tabular}
\column{0.3\textwidth}
\begin{tabular}{@{}cl@{}}
\toprule
\multicolumn{2}{c}{Marker}\\
\midrule
\kd{o} & lingkaran\\ \kd{s} & kotak\\ \kd{\textasciicircum} & segitiga\\
\kd{d} & belah ketupat\\ \kd{p} & bintang 5\\ \kd{*} & bintang\\
\kd{+} & plus\\ \kd{x} & silang\\ \kd{.} & titik\\
\bottomrule
\end{tabular}
\end{columns}

\medskip
Ketiganya digabung dalam satu teks, \alert{urutan bebas}: \kd{plot(x, y, 'r-{}-o')} $\rightarrow$ merah, putus-putus, marker lingkaran. Jika marker ditulis tanpa gaya garis (\kd{'o'}), hanya titik yang digambar.
\end{frame}

\begin{frame}[fragile]{Jenis gaya garis (\emph{line style})}
\kiri{211pt}
\konsep Gaya garis membedakan variabel \alert{tanpa bergantung pada warna}, penting untuk cetakan hitam-putih atau pembaca buta warna.
\begin{lstlisting}
x = 1:10;
hold on
plot(x, x,   '-')    % lurus (bawaan)
plot(x, x+2, '--')   % putus-putus
plot(x, x+4, ':')    % titik-titik
plot(x, x+6, '-.')   % garis-titik
hold off

% atau lewat properti:
% plot(x, x, 'LineStyle','--')
\end{lstlisting}
\kd{'LineStyle','none'} menyembunyikan garis sehingga hanya marker yang tampak.
\kanan{211pt}{f16_linestyle}
\end{frame}

\begin{frame}[fragile]{Jenis marker}
\kiri{211pt}
\konsep Marker menandai \alert{lokasi data yang sebenarnya}; garis di antaranya hanyalah penghubung.
\begin{lstlisting}
x = 1:8;
hold on
plot(x, x,    'o')   % lingkaran
plot(x, x+2,  's')   % kotak
plot(x, x+4,  '^')   % segitiga
plot(x, x+6,  'd')   % belah ketupat
plot(x, x+8,  'p')   % bintang 5
plot(x, x+10, '*')   % bintang
hold off
\end{lstlisting}
Karena hanya simbol marker yang ditulis, garis tidak digambar. Gabungkan dengan gaya garis dan warna untuk garis + marker, misalnya \kd{'r-{}-s'}.
\kanan{211pt}{f17_marker_jenis}
\end{frame}

\begin{frame}[fragile]{Mengatur marker pada garis}
\kiri{182pt}
\konsep Untuk data tahunan yang jumlahnya sedikit, garis + marker membuat setiap pengamatan terlihat.
\begin{lstlisting}
hold on
plot(tahun, growth,  '-o', ...
     'MarkerFaceColor','#0072BD')
plot(tahun, inflasi, '--s', ...
     'MarkerSize',8)
plot(tahun, bunga,   ':^', ...
     'MarkerSize',10, ...
     'MarkerFaceColor','y')
hold off
\end{lstlisting}
\kd{MarkerSize}: ukuran (bawaan 6). \kd{MarkerFaceColor}: warna isi (bawaan \kd{'none'} = kosong). \kd{MarkerEdgeColor}: warna tepi.
\kanan{182pt}{f12_marker}
\end{frame}

\begin{frame}[fragile]{Membuat grid}
\kiri{207pt}
\konsep Grid membantu mata ``menarik garis'' dari titik data ke angka pada sumbu.
\begin{lstlisting}
plot(tahun, growth)
grid on      % grid utama
grid minor   % + grid minor
% grid off   % matikan
\end{lstlisting}
Grid utama mengikuti posisi \emph{tick}; grid minor membagi lagi di antaranya. Gunakan grid minor seperlunya: terlalu banyak garis justru mengganggu.
\kanan{207pt}{f07_grid}
\end{frame}

\begin{frame}[fragile]{Rentang dan \emph{tick} sumbu}
\kiri{182pt}
\konsep MATLAB memilih rentang sumbu otomatis. Atur sendiri jika ingin menyamakan skala antar-grafik atau memberi ruang bagi legenda.
\begin{lstlisting}
plot(tahun, growth)
xlim([2018 2024])   % batas sumbu x
ylim([-5 10])       % batas sumbu y
xticks(2018:2024)   % posisi angka
grid on
% x dan y sekaligus:
% axis([2018 2024 -5 10])
% kembali otomatis: ylim auto
\end{lstlisting}
\kd{xlim}/\kd{ylim} menerima \kd{[min max]}. \kd{xticks}/\kd{yticks} menentukan di mana angka ditulis; di sini agar setiap tahun muncul tanpa angka desimal.
\kanan{182pt}{f11_sumbu}
\end{frame}

\begin{frame}[fragile]{Skala logaritmik}
\small
\konsep Pada skala log, jarak yang sama berarti \alert{rasio} yang sama. Pertumbuhan dengan laju konstan (eksponensial) tampak sebagai garis lurus.
\begin{columns}[T,onlytextwidth]
\column{\dimexpr\textwidth-219pt-1.2em\relax}
\begin{lstlisting}
t   = 0:30;
pdb = 100 * 1.05.^t;   % tumbuh 5%/th

plot(t, pdb)           % kiri: linear
semilogy(t, pdb)       % kanan: log

% setara dengan:
% plot(t, pdb)
% set(gca, 'YScale', 'log')
\end{lstlisting}
Ada juga \kd{semilogx} dan \kd{loglog}. Skala log hanya untuk nilai \alert{positif}; nilai nol atau negatif tidak digambar.
\column{219pt}
\includegraphics[width=219pt]{figs/f26_log}
\end{columns}
\end{frame}

\begin{frame}[fragile]{Diagram dua sumbu $y$ (kiri--kanan)}
\kiri{182pt}
\konsep Jika dua variabel punya skala atau satuan berbeda, masing-masing diberi sumbu $y$ sendiri.
\begin{lstlisting}
yyaxis left
plot(tahun, growth, '-o')
ylabel('Growth (%)')

yyaxis right
plot(tahun, bunga, '-s')
ylabel('Suku bunga (%)')

xlabel('Tahun')
title('Growth dan Suku Bunga')
\end{lstlisting}
\kd{yyaxis left}/\kd{right} memilih sumbu aktif; perintah plot berikutnya masuk ke sumbu itu. Hati-hati: rentang kedua sumbu bebas diatur, sehingga ``persilangan'' garis bisa \alert{menyesatkan}.
\kanan{182pt}{f14_yyaxis}
\end{frame}

\begin{frame}[fragile]{Mengubah properti setelah plot dibuat (\emph{handle})}
\small
\konsep Setiap objek grafik bisa ``dipegang'' lewat sebuah variabel (\emph{handle}). Propertinya bisa dibaca dan diubah kapan saja, tanpa menggambar ulang.
\begin{columns}[T,onlytextwidth]
\column{0.5\textwidth}
\begin{lstlisting}
h = plot(tahun, growth);   % h = objek Line
h.LineWidth = 2;           % ubah tebal
h.Color     = '#D95319';   % ubah warna
h.Marker    = 'o';         % tambah marker

ax = gca;                  % axes aktif
ax.FontSize = 12;          % ukuran huruf
ax.YLim     = [-5 10];     % = ylim([-5 10])

get(h)                     % lihat semua properti
set(h, 'LineStyle', '--')  % cara lama
\end{lstlisting}
\column{0.46\textwidth}
{\footnotesize
\begin{tabular}{@{}ll@{}}
\toprule
Fungsi & Mengembalikan\\
\midrule
\kd{gcf} & figure aktif (\emph{get current figure})\\
\kd{gca} & axes aktif (\emph{get current axes})\\
\kd{h = plot(\ldots)} & objek garis yang baru dibuat\\
\bottomrule
\end{tabular}}

\medskip
Semua nama-nilai yang sudah dipelajari (\kd{'LineWidth'}, \kd{'Color'}, \ldots) sebenarnya adalah \alert{properti} objek ini. Notasi titik (\kd{h.LineWidth}) berlaku sejak R2014b.
\end{columns}
\end{frame}

% =========================================================================
\section{Beberapa grafik dan menyimpan hasil}

\begin{frame}[fragile]{Beberapa grafik dalam satu figure}
\kiri{182pt}
\konsep Figure dibagi menjadi kisi berisi beberapa axes; \kd{nexttile} pindah ke kotak berikutnya.
\begin{lstlisting}
tiledlayout(2, 2)   % 2 baris x 2 kolom

nexttile
plot(tahun, growth);  title('Growth')
nexttile
plot(tahun, inflasi); title('Inflasi')
nexttile
bar(tahun, bunga);    title('Suku bunga')
nexttile
scatter(inflasi, growth, 'filled')
title('Inflasi vs growth')
\end{lstlisting}
\kd{tiledlayout} ada sejak R2019b. Cara lama: \kd{subplot(2,2,k)}, $k$ = 1--4 dihitung per baris. Setiap kotak adalah axes tersendiri; \kd{title}, \kd{xlabel}, dsb.\ berlaku untuk kotak yang aktif.
\kanan{182pt}{f28_tiled}
\end{frame}

\begin{frame}[fragile]{Menyimpan grafik ke file}
\small
\konsep Grafik di layar perlu disimpan ke file agar bisa dipakai di laporan atau slide. Format dipilih sesuai tujuan.
\begin{lstlisting}[basicstyle=\ttfamily\footnotesize]
exportgraphics(gcf, 'growth.png', 'Resolution', 300)       % raster 300 dpi
exportgraphics(gcf, 'growth.pdf', 'ContentType', 'vector') % vektor
savefig(gcf, 'growth.fig')   % format MATLAB, bisa diedit lagi
saveas(gcf, 'growth.png')    % cara lama
\end{lstlisting}
{\footnotesize
\begin{tabular}{@{}p{0.15\textwidth}p{0.42\textwidth}p{0.36\textwidth}@{}}
\toprule
Format & Sifat & Cocok untuk\\
\midrule
PNG & Raster (piksel); pecah jika diperbesar & Word, PowerPoint, web\\
PDF / EPS & Vektor; tetap tajam di ukuran berapa pun & \LaTeX, publikasi\\
FIG & Format MATLAB; objek grafik tetap utuh & Mengedit grafik di kemudian hari\\
\bottomrule
\end{tabular}}

\medskip
\kd{exportgraphics} (R2020a+) otomatis memotong ruang putih berlebih. Ganti \kd{gcf} dengan \kd{gca} untuk menyimpan satu axes saja.
\end{frame}

% =========================================================================
\section{Penutup}

\begin{frame}[fragile]{Contoh gabungan: dari data sampai file}
\begin{columns}[T,onlytextwidth]
\column{\dimexpr\textwidth-154pt-1.2em\relax}
\begin{lstlisting}
figure                          % (2)
hold on                         % (3)
plot(tahun, growth,  '-o', 'LineWidth',2, ...
     'Color','#0072BD', 'DisplayName','Growth')
plot(tahun, inflasi, '--s','LineWidth',2, ...
     'Color','#D95319', 'DisplayName','Inflasi')
plot(tahun, bunga,   ':^', 'LineWidth',2, ...
     'Color','#77AC30', 'DisplayName','Suku bunga')
hold off
grid on                         % (4)
xlim([2018 2024]); xticks(2018:2024)
title('Indikator Ekonomi'); subtitle('Data ilustrasi')
xlabel('Tahun'); ylabel('Persen (%)')
legend('Location','southoutside', ...
       'Orientation','horizontal')
% (5) simpan
exportgraphics(gcf,'indikator.png','Resolution',300)
\end{lstlisting}
\column{154pt}\small
\includegraphics[width=154pt]{figs/f15_gabungan}\par\medskip
Nomor (2)--(5) pada komentar mengikuti \alert{alur kerja} di awal materi; langkah (1) adalah data di slide pembuka.
\end{columns}
\end{frame}

\begin{frame}{Kesalahan yang sering terjadi}
\footnotesize
\begin{tabular}{@{}>{\raggedright\arraybackslash}p{0.32\textwidth}>{\raggedright\arraybackslash}p{0.27\textwidth}>{\raggedright\arraybackslash}p{0.33\textwidth}@{}}
\toprule
Gejala & Penyebab & Solusi\\
\midrule
Error \kd{Vectors must be the same length.} & Panjang $x \neq y$ & Cek \kd{numel(x)}, \kd{numel(y)}\\
Garis yang digambar sebelumnya hilang & Lupa \kd{hold on} & Pasang \kd{hold on} sebelum plot kedua\\
Grafik lama tertimpa grafik baru & Tidak membuka jendela baru & Panggil \kd{figure} dulu\\
Nama di legenda tertukar & Urutan \kd{legend} $\neq$ urutan plot & Pakai \kd{'DisplayName'}\\
Judul \kd{GDP\_growth} tampil sebagai subskrip & Interpreter \TeX & \kd{'Interpreter','none'}\\
Error saat menghitung \kd{y = x\textasciicircum 2} & Operasi matriks, bukan per elemen & Pakai \kd{x.\textasciicircum 2}, \kd{.*}, \kd{./}\\
Garis zig-zag bolak-balik & Data \kd{x} tidak terurut & \kd{[x,i] = sort(x); y = y(i);}\\
Garis tidak terlihat sama sekali & Data hanya satu titik, atau ada \kd{NaN} & Tambah marker (\kd{'o'}), cek \kd{isnan}\\
\bottomrule
\end{tabular}
\end{frame}

\begin{frame}{Ringkasan perintah}
\scriptsize
\begin{columns}[T,onlytextwidth]
\column{0.48\textwidth}
\begin{tabular}{@{}ll@{}}
\toprule
Tujuan & Perintah\\
\midrule
Plot garis   & \kd{plot(x,y)}, \kd{plot(x,Y)}\\
Batang       & \kd{bar(x,y)}, \kd{'stacked'}\\
Sebar        & \kd{scatter(x,y,ukuran)}\\
Jenis lain   & \kd{histogram}, \kd{area}, \kd{stairs}, \kd{pie}\\
Garis tren   & \kd{polyfit}, \kd{polyval}\\
Jendela      & \kd{figure}, \kd{clf}, \kd{close all}\\
Tambah objek & \kd{hold on}, \kd{hold off}\\
Judul        & \kd{title}, \kd{subtitle}\\
Label sumbu  & \kd{xlabel}, \kd{ylabel}\\
Legenda      & \kd{legend}, \kd{'DisplayName'}, \kd{'Location'}\\
Anotasi      & \kd{xline}, \kd{yline}, \kd{text}\\
\bottomrule
\end{tabular}
\column{0.48\textwidth}
\begin{tabular}{@{}ll@{}}
\toprule
Tujuan & Perintah\\
\midrule
Gaya singkat   & \kd{'r-{}-o'} (\emph{LineSpec})\\
Tebal / warna  & \kd{'LineWidth'}, \kd{'Color'}\\
Marker         & \kd{'Marker'}, \kd{'MarkerSize'}\\
Grid           & \kd{grid on}, \kd{grid minor}\\
Rentang sumbu  & \kd{xlim}, \kd{ylim}, \kd{axis}\\
Posisi angka   & \kd{xticks}, \kd{yticks}\\
Skala log      & \kd{semilogy}, \kd{'YScale','log'}\\
Dua sumbu $y$  & \kd{yyaxis left}, \kd{yyaxis right}\\
Properti objek & \kd{h = plot(\ldots)}, \kd{gca}, \kd{gcf}\\
Banyak grafik  & \kd{tiledlayout}, \kd{nexttile}\\
Simpan         & \kd{exportgraphics}, \kd{savefig}\\
\bottomrule
\end{tabular}
\end{columns}
\end{frame}

\begin{frame}{Latihan}
Gunakan data \kd{tahun}, \kd{growth}, \kd{inflasi}, dan \kd{bunga} dari slide pembuka.
\begin{enumerate}
  \item \textbf{Dasar.} Buat grafik inflasi terhadap tahun dengan garis \alert{merah putus-putus} dan marker kotak. Beri judul, label sumbu lengkap dengan satuan, dan grid.
  \item \textbf{Perbandingan.} Buat diagram batang berkelompok untuk growth dan inflasi. Letakkan legenda di bawah grafik secara horizontal. Pada tahun berapa inflasi lebih tinggi daripada growth?
  \item \textbf{Hubungan.} Buat scatterplot suku bunga ($x$) terhadap inflasi ($y$) beserta garis tren. Ulangi tanpa tahun 2020. Apakah kemiringan garisnya banyak berubah?
  \item \textbf{Laporan.} Gabungkan grafik dari soal 1 dan 3 dalam satu figure (\kd{tiledlayout(1,2)}), lalu simpan sebagai PNG 300 dpi dan sebagai PDF vektor.
\end{enumerate}
\smallskip
{\small Petunjuk soal 3: tahun 2020 adalah elemen ke-3; buang dengan \kd{idx = tahun \textasciitilde= 2020;} lalu pakai \kd{bunga(idx)} dan \kd{inflasi(idx)}.}
\end{frame}

\end{document}
